\section{Gasping for Air}

\begin{quotex}
A disciple asked his spiritual teacher, while they were bathing in the river, when he would finally be able to realise Brahman. The teacher, instead of answering, pushed his head under water and held him down until, feeling himself drowning, the student freed himself and re-emerged. The teacher then explained: ``When the desire in you to realise Brahman becomes as intense and deep as how you were just driven to reassert your physical life---only then will you achieve satisfaction.'' 
\flright{\textsc{Julius Evola}, \emph{Essays on Magical Idealism}\footnote{\url{https://gornahoor.net/?p=20}}}
\end{quotex}


A recent flare up of bronchitis left me gasping for air for days at a time and unable to write. It was a reminder of death (that same weekend a famous architect was admitted to a hospital in Miami with bronchitis where she died) and stoked the desire to ``realise Brahman''.

\paragraph{Like a Herd of Cud Chewers}


I've been a fan of nature shows ever since I saw Mutual of Omaha's Wild Kingdom\footnote{\url{http://www.wildkingdom.com/}}. I especially relate to big cats and was fascinated by their coordinated attacks on herd animals. What struck me the most was how the cud chewers would quickly return to their normal activities after one of their own had been picked off. As a naïve youth, I thought that humans would never act that way.

However, in current times after an attack on a population, the government encourages everyone to get back to ``normal'', even to start ``shopping'' again. Flowers are laid, candles are lit, and people reassure each other about their ``resiliency''. Anything to avoid facing up to the real solution. That is because those who kill the soul are more dangerous than those who kill the body.

\paragraph{First Order Logic}


Years ago at University, I took a class on first order logic (and Zermelo-Fraenkel set theory) in the philosophy department. First order logic\footnote{\url{https://en.wikipedia.org/wiki/First-order_logic}} is about statements with predicates and quantifiers, such as the universal and existential quantifiers. In other words, it is about statements such as ``All X are P'' or ``Some X are P'', and so on, pretty much covering most everyday discourse.

I recall that about two dozen signed up for the class, but only three of us finished. I knew several of the drop outs: they were majoring in political science, pre-law, even philosophy. In short, those who most needed it. Today on the news shows, the talking heads always struggle with justifications such as ``I don't mean that all X are P'' and so on, before they can get to the main point. Needless, the X's in the audience object to being called P. Perhaps if every university graduate had to understand rudimentary logic, discourses would be more fruitful.

\paragraph{Heidegger ... with friends like these}

No one who knows me personally reads Gornahoor; I suppose you should draw your own conclusions. My sister thinks I am a shaman, who should perhaps be avoided. Others, trying to be helpful, send me links to various spiritual conferences, visiting gurus, and the like. I wonder why they never send them my link.

The end result is that I've accumulated a new type of 21\textsuperscript{st} century relationship without a name, so I'll use the term ``Internet friend''. Usually, they are just for a season, or those with a specific common interest. One such ``friend'', whom I've ``known'' off and on for years, recently wrote to me about Heidegger and Tradition, specifically Evola ... In particular, Heidegger is the cat's meow among the ``alt right''.

It brings up some questions: does the language of Tradition, with Sanskrit terms, obscure more than reveal? Can Heidegger's language be used to express the same ideas? So, instead of ``karma'', we can speak of ``facticity'', instead of dharma, we talk about projects. And for human nature, we can talk about man's mode of being in the world.

Of course, we would supplement Heidegger's phenomenological analysis with descriptions of the inner life of the soul. I think the biggest benefit would be to abandon any talk of the ``Kali Yuga''. It is too easily portrayed as some pseudo-mechanical process that operates independently of any consciousness. Man becomes, in that view, a passive spectator, rather than the active agent of change.

Instead, we can talk about nihilism a la Heidegger. This is clearer, when nihilism is understood as the complete loss of Tradition. As we've pointed out on these pages, this is the age of the rule by scientists, or technocrats, and anthropoidal mammals. Those who have read carefully will begin to notice more similarities.

\paragraph{The Zen Master Dogen}


H/T to Aryayana for drawing our attention to The Life of Zen Master Dogen\footnote{\url{https://youtu.be/8s5Uz8XSc-Y}}. No car chases or CGI monsters, if that is your preference. Rather, it provides a rather convincing portrayal of Tradition in 13\textsuperscript{th} century Japan. There is the ongoing relationship between the spiritual authority and temporal power. Dogen's recovery of true Buddhism involves the discovering of the ``Buddha within''. He illustrates that with the reflection of the full moon in a pond. The surface of the water represents the Self; when it is disturbed, the moon is not reflected very well. However, when the water is still, the moon is reflected perfectly. Ignorance and vice disturb the Self so that the image of the Buddha cannot be found. Hence, the life of virtue is necessary, which has 8 steps, including meditation. Sitting zazen requires a good posture and the elimination of all fidgeting.

At that level, a man of one Tradition can engage a man of another tradition. They can discuss meditation practice, techniques to still the mind, and aids to overcome vices like gluttony, lust, pride and so on. That is real life. In existentialist terms, we can say it involves the recovery of being rather than being engaged in debates about theories.

That is why we object to meaningless assertions such as ``X is a valid tradition'', or a ``sect of Old Believers has maintained tradition''. That is good first order logic, but not good for spiritual practice. The nominal attachment to a Tradition leads to nothing, just as Dogen showed that inauthentic forms of Buddhism do not lead to enlightenment.

\paragraph{Imaginal and Conceptual Thinking}


The distinction between imaginal and conceptual thinking is difficult for most people to make. \textbf{Rene Guenon} has claimed that those who cannot think outside of time are also incapable of proper metaphysical thinking. Nevertheless, since knowledge begins in the senses, symbols and imagery are used to convey metaphysical teachings; these involve visualizations in space and time.

In particular, postmortem states are described in terms of images rather than in concepts. Islam is famous, or infamous, for that, with its very sensuous images of ``heaven''; of course, there is an esoteric interpretation of such texts. On the other hand, most Americans look forward in the afterlife, not to the beatific vision of God, but rather to being reunited with their favorite Chihuahua.

\paragraph{Mortal Sin}


Mortal sin is a term that is easily tossed around. Why, it is asked, does a finite sin have infinite consequences? Mortal sin just means a second death, the death of the soul. So it is just as permanent as the first death of the body.

If you are reading text messages while crossing the street, and get run over, isn't the penalty too severe for the offense? It doesn't matter, does it? If a man beats his son or molests his daughter, and they despise him for life, does a finite act deserve a lifelong punishment?

To be clear, a mortal sin requires three elements:

\begin{itemize}


\item A serious matter

\item Knowledge that it is serious

\item The act is deliberate
\end{itemize}


In short, only the man who knows and the man who is free can commit such sins. He would look at that prospect with total horror. That is not unlike the young chef in the Dogen movie who was horrified to discover his lust for a pretty girl.

\paragraph{Provocations}


Back to Dogen: how can we in the west overcome ignorance and vice? It turns out that the two are intimately related. \textbf{St Thomas Aquinas} explains why most men remain in ignorance:

\begin{quotex}
in some the reason is perverted by passion, or evil habit, or an evil disposition of nature
\end{quotex}


So the solution to the lack of gnosis is not to read another book or study a philosopher on-line; these are pointless without restraining passions and overcoming evil habits. The Fathers described eight provocations that keep us in bondage: gluttony, lust, greed, anger, depression, boredom, vanity, and pride.

The effects of these provocations need to be observed in consciousness through watchfulness, and the kept from dominating. The phenomenology of provocations will be known both to Christians and to Zen monks.

A good contemporary resource is \emph{Sin Revisited} by \textbf{Solange Hertz}, who describes them rather well. She includes the esoteric interpretation of some Old Testament texts, which need to be understood as describing spiritual warfare.


\flrightit{Posted on 2016-04-11 by Cologero}

\begin{center}* * *\end{center}

\begin{footnotesize}
\begin{sffamily}
\texttt{mtpollack on 2016-04-11 at 13:49 said:}

\emph{``What struck me the most was how the cud chewers would quickly return to their normal activities after one of their own had been picked off. As a naïve youth, I thought that humans would never act that way.}

It's an all-you-can eat Kunda-buffet.

\hfill

\texttt{Aryayana on 2016-04-12 at 14:54 said:}

I am glad that you appreciated the Zen film just as much as I did, Cologero, and your reflections are very pertinent.

Alain de Benoist is big on Heidegger. I tried to read his `Being and Time' back in the day, but decoding his extremely dense and abstruse personal philosophical construct didn't seem to be worth the time, when there are so many truly timeless works to read instead. One odd thing I have noticed is that certain people attempt to `synthesize' (perhaps rather syncretize) the Traditional thought of an Evola or a Guenon with other, non-traditional philosophers like Heidegger even before having read Guenon's entire oeuvre, not to mention the works of Schuon, Coomaraswamy et al (and even less than that, having established oneself firmly in a particular stream of spiritual doctrine and praxis).

On another note, I have to disagree about dropping terminological reference to the Kali-yuga.

Completely regardless of whether or not it is mentioned, it remains the truth about what is happening; therefore, it serves to see the entire situation in light of the greater picture, to orient oneself realistically before deciding which course of action is most appropriate in one's own case. For those whose mission (rightly or wrongly) is to work to change the world through external projects, thoughts about the inevitability of further collective degeneration may be demoralizing. For those of a contemplative nature, however, this awareness only reinforces a sense of the urgency of transcending this world by means of the esoteric Path. It definitely eliminates the temptation to abandon an exclusive spiritual quest well befitting the caste of sages in favour of all sorts of worldly projects aimed at ``saving'' the world on a collective level, certainly doomed to fail. Heroic (in a tragic sense) as such a mission might be, it practically boils down to a struggle of individuals versus the entire world, amounting to spending for nothing a life that could have been devoted to intensive spiritual practice (the far-reaching, magnetic polar influence of which should not be underestimated) and transcendence at the one and only level where this is still possible in our time: that of the individual himself. That the world-situation depends upon consciousness is a given, but mankind (with the exception of a few) has already made its choice, and the mighty forces (born of far more than our own consciousness) that have been unleashed must be allowed to run their course, until this manifestation is exhausted in final ruin. (However central one thinks one's personal role may be, the world goes on even after the death of the most extraordinary men, many of whose efforts were unable to turn a tide that was too mighty for them: think of what became of Codreanu's mission, for example.) The fate of the world is in God's hands, and on a collective level that fate looks bleak, but for the individual, there is available a path that may defeat the enemy once and for all in our own souls, no matter what becomes of the world at large. To completely confront, without fear, the likely possibility that the world is in fact utterly lost to darkness, with no return to a true sacred order in sight, and to calmly and serenely accept it, consequently letting go of all worldly hopes and interests, may do wonders for the sense of absolute detachment that can greatly benefit the Great Work. Thus the wise man doesn't become a passive spectator, but an active agent in overcoming the Dark Age within himself, rather than running around imagining himself to be able to change the way of the rest of the world.

The present world is comparable to a collective train that is headed for the abyss; a train whose movement cannot be stopped, nor its course be changed, since it is controlled by forces far outweighing the power of those individuals opposed to it, but that can be exited or renounced by individual passengers if they only dare to jump off. Even though the impetus of the train's movement towards the abyss is at this point much too forceful to be stopped by any human effort, we shouldn't for that reason personally submit to its determined course by thinking that the inevitable equals the Good, instead of choosing the exit, like a silent warrior in the night, so as to keep the fire of Truth alive in one more heart, even to the end. On the historical plane, this world will be destroyed sooner or later (an obvious fact which reminds us that the essential Quest of man is not tied up with the external salvation of the world, which is ultimately outside man's control) a demonstration of the illusory and vain nature of all human power and control when faced with the lightning-flash of the Absolute which destroys, reintegrates and remanifests. Then all the efforts of the Enemy who thought they had finally managed to subvert humanity entirely will once again be nothing, for their power was always nothing in view of the Ultimate, and all they really achieved was to spend their precious human existences striving so as to fall into lower states of being instead of rising to higher ones. This is the tragedy of relative power gained from darkness, so to speak. It comes at a high cost and is impermanent, as it only pertains to the most transient and limited domain. Everything that has a beginning, has an end, and to think that human self-will can change that in any other way than through its reintegration into the Divine Source is deluded. And this is where the doctrine of cycles or yugas is entirely correct. The consciousness of a being will certainly determine the `yugic' conditions in which he manifests following the disintegration of his current form (from which his external environment---and thus the Dark Age---is indeed inseparable), but while situated in a cyclic manifestation shared with countless other beings, he cannot single-handedly change those external conditions as if he were Almighty God. The conscious, individual agents who have worked with the utmost determination to create the nihilistic chaos we have today, may certainly believe that this is the pure result of their own ``power'' and ``will'', but the truth is that they are nothing, they are mere pawns to forces out of their control, and no one is more trapped than they are. They have exploited the possibilities that the cyclic conditions already allowed for, and nothing is easier than that, if one doesn't see a problem in causing the damnation of one's soul. For those who are not determined to make sure that the death of the cycle will be accompanied by one's own spiritual death, however, one of the few remaining options is to leave the world to its fate and invest all their forces in spiritual reintegration, a path that must be the opposite pole (``... my Kingdom is not of this world'') to those who have contrarily chosen to make this dark world-age their own kingdom, changing themselves in conforming thereto ... as the time before the line was crossed allowing for the possibility of temporary (and increasingly short-lived) returns to higher orders has long since passed. Maybe this is a truth that is better kept silent, though, since not all people find themselves motivated by it as do I. It helps to remember that all Principles that truly matter, including the Spiritual Archetype of Man in the Divine Mind, outlast their contingent manifestations or emanations in this limited and particular world-cycle, a cycle whose duration cannot be prolonged indefinitely on the initiative of men.

\hfill

\texttt{undrownpip on 2016-04-17 at 20:02 said:}

I kinda disagree about the Kali Yuga --- yes it's yielding to some vast mechanical tide, but it's a self-conscious yielding, a giving over to what is inevitable. And what is inevitable is inevitable (the ``ride the tiger'' crowd thinks) because they are telling the world exactly what to do to save itself, and nobody is listening. At that point, what can you do but accept destruction?

\hfill

\texttt{Mark Citadel on 2016-04-18 at 18:39 said:}

I don't think I could ever jettison Kali Yuga from my vocabulary. haha. It's just become too ingrained. But I do think you're onto something concerning Heidegger there. I have been wanting for a while to read Dugin's treatment of Heidegger which was relatively recent and no doubt tries to square him with eastern thought.

Sorry to hear you have been unwell! Please don't die on us.

\hfill

\texttt{Nilakantha on 2016-04-30 at 02:25 said:}

I've always had a hard time thinking of any of the Protestant Buddhisms of the Kamakura as Traditional in any important sense. Even Dogen's Zen is under the cloud of subitism and is, therefore, heterodox.

\hfill

\texttt{Cologero on 2016-04-30 at 12:11 said:}

Nilakantha, I can't comment on the merits or demerits of your point. However, that uncertainty was one of the main reasons why I decided to pass on Buddhism: I came to the conclusion that I could not readily distinguish between the authentic and the heterodox in an alien tradition.

\hfill

\texttt{Nilakantha on 2016-04-30 at 14:48 said:}

Cologero,

Decades ago, when I left the Orthodox Church, I knew that I would need to find a spiritual home in an orthodox religion, one that had a well articulated metaphysical doctrine, intelligible morality and traditional spiritual practices. Guénon, Coomaraswamy and Plotinus helped keep me honest in my search. I also took the advice of Agehananda Bharati (of all people!).

When Bharati was teaching at Benares Hindu University in the 60s and 70s, he would ask the young western hippies who wanted to study Buddhism with him if they had done their homework, i.e. had they learned any of the languages of traditional Buddhism: Pali, Sanskrit, Classical Chinese or Tibetan. If they hadn't, he wouldn't take them as pupils, saying that these were the only languages that had coherent terminology for explaining Buddhism. I agree with Bharati here, so I spent six years in graduate school studying Sanskrit, Pali and classical Indian thought.

I only bring this up to show that I agree with your statement that one `` can not readily distinguish between the authentic and the heterodox in an alien tradition.'' If for some reason you need to leave your ancestral religion, it is incumbent upon you to become as much a native of your new faith as possible.

\hfill

\texttt{Cologero on 2016-05-11 at 05:47 said:}

Mark, it is always a \textit{beautiful day to die}\footnote{\url{https://www.youtube.com/watch?v=r2ILyyWh3sY}}. But there are some loose ends:

\begin{itemize}

\item Who will keep the web site up?

\item What will happen to my nice library? I would hate to see it carted off to the used book store.

\item Who will pray for my soul?
\end{itemize}


\end{sffamily}
\end{footnotesize}

